%
%Map ID hợp nhất 4 mức độ.
%
%Nếu là ID5 sẽ theo định đạng: %[Tham số 1 Tham số 2 Tham số 3 Tham số 4 Tham số 5]. Ví dụ: %[1D2B3]
%Nếu là ID6 sẽ theo định đạng: %[Tham số 1 Tham số 2 Tham số 3 Tham số 4 Tham số 5 - Tham số 6]. Ví dụ: %[1D2B3-1]
%Có thể thay đổi nội dung của mức độ nhưng vị trí của mức độ trong ID không đổi (thông số thứ 3 từ trái qua).
%Cú pháp của thông số: [Giá trị] Mô tả thông số. Ví dụ: [0] Lớp 10 thì 0 là giá trị để lưu vào ID.
%
%Chú ý: Các dấu gạch ngang không được thay đổi.
%
%Cấu hình tên các tham số
%
Tên tham số 1: Lớp
Tên tham số 2: Môn
Tên tham số 3: Chương
Tên tham số 4: Mức độ
Tên tham số 5: Bài
Tên tham số 6: Dạng
%
%Cấu hình chi tiết ID
%
%%Cấu hình mức độ dùng chung.
[N] Nhận biết
[H] Thông hiểu
[V] Vận dụng
[C] Vận dụng cao
%
%Cấu hình nội dung
%
-[0] Lớp 10
----[D] Đại số/Thống kê 10
-------[1] Mệnh đề. Tập hợp
----------[1] Mệnh đề
-------------[1] Xác định mệnh đề, mệnh đề chứa biến
-------------[2] Tính đúng-sai của mệnh đề
-------------[3] Phủ định của một mệnh đề
-------------[4] Mệnh đề kéo theo, mệnh đề đảo, mệnh đề tương đương
-------------[5] Mệnh đề với mọi, tồn tại (và phủ định chúng)
-------------[6] Áp dụng mệnh đề vào suy luận có lí
----------[2] Tập hợp
-------------[1] Tập hợp và phần tử của tập hợp
-------------[2] Tập hợp con. Hai tập hợp bằng nhau
-------------[3] Ký hiệu khoảng, đoạn, nửa khoảng
----------[3] Các phép toán tập hợp
-------------[1] Giao và hợp của hai tập hợp (rời rạc)
-------------[2] Hiệu và phần bù của hai tập hợp (rời rạc)
-------------[3] Giao và hợp (dùng đoạn, khoảng)
-------------[4] Hiệu và phần bù (dùng đoạn, khoảng)
-------------[5] Toán thực tế ứng dụng của tập hợp
-------[2] BPT và hệ BPT bậc nhất hai ẩn
----------[1] Bất phương trình bậc nhất hai ẩn
-------------[1] Các khái niệm về BPT bậc nhất hai ẩn
-------------[2] Miền nghiệm của BPT bậc nhất hai ẩn
-------------[3] Toán thực tế về BPT bậc nhất hai ẩn
----------[2] Hệ bất phương trình bậc nhất hai ẩn
-------------[1] Các khái niệm về Hệ BPT bậc nhất hai ẩn
-------------[2] Miền nghiệm của Hệ BPT bậc nhất hai ẩn
-------------[3] Toán thực tế về Hệ BPT bậc nhất hai ẩn
-------[3] Hàm số bậc hai và đồ thị
----------[1] Hàm số và đồ thị
-------------[1] Xác định một hàm số
-------------[2] Tập xác định của hàm số
-------------[3] Giá trị của hàm số
-------------[4] Đồ thị của hàm số
-------------[5] Tính đồng biến, nghịch biến
-------------[6] Tính chẵn, lẻ
-------------[7] Toán thực tế về hàm số
----------[2] Hàm số bậc hai
-------------[1] Xác định hàm số bậc hai
-------------[2] Bảng biến thiên, tính đơn điệu, max, min
-------------[3] Đồ thị của hàm số bậc hai
-------------[4] Bài toán về sự tương giao
-------------[5] Hàm số chứa dấu giá trị tuyệt đối
-------------[6] Toán thực tế ứng dụng hàm số bậc hai
-------[6] Thống kê
----------[1] Số gần đúng. Sai số
-------------[1] Tính và ước lượng sai số tuyệt đối, tương đối
-------------[2] Tính và xác định độ chính xác của kết quả
-------------[3] Quy tròn số gần đúng
-------------[4] Viết số gần đúng cho số đúng biết độ chính xác
----------[2] Mô tả và biểu diễn dữ liệu trên các bảng và biểu đồ
-------------[1] Đọc và phân tích thông tin trên bảng số liệu
-------------[2] Đọc và phân tích thông tin trên biểu đồ
-------------[3] Số liệu bất thường trên bảng số liệu
-------------[4] Số liệu bất thường trên biểu đồ
----------[3] Các số đặc trưng đo xu thế trung tâm của mẫu số liệu
-------------[1] Câu hỏi lý thuyết
-------------[2] Số trung bình cộng
-------------[3] Số trung vị
-------------[4] Tứ phân vị
-------------[5] Mốt
----------[4] Các số đặc trưng đo mức độ phân tán của mẫu số liệu
-------------[1] Câu hỏi lý thuyết
-------------[2] Khoảng biến thiên, khoảng tứ phân vị
-------------[3] Giá trị bất thường của mẫu số liệu
-------------[4] Phương sai, độ lệch chuẩn của mẫu số liệu
-------[7] Bất phương trình bậc 2 một ẩn
----------[1] Dấu của tam thức bậc 2
-------------[1] Xác định tam thức bậc 2
-------------[2] Dấu của tam thức bậc 2 và ứng dụng
-------------[3] Bài toán xét dấu biết BXD, đồ thị
-------------[4] Xét dấu biểu thức dạng tích, thương
-------------[5] Toán thực tế ứng dụng dấu tam thức bậc 2
----------[2] Giải bất phương trình bậc 2 một ẩn
-------------[1] Bất phương trình bậc 2 và ứng dụng
-------------[2] Giải bất phương trình bậc hai biết BXD, đồ thị
-------------[3] Bất phương trình dạng tích, thương
-------------[4] Hệ bất phương trình BPT bậc 2
-------------[5] Bất phương trình chứa căn, | · |
-------------[6] Bài toán có tham số m
-------------[7] Toán thực tế ứng dụng bất phương trình bậc 2 một ẩn
----------[3] Phương trình quy về phương trình bậc hai
-------------[1] Phương trình căn √︀f(x) = √︀g(x) và mở rộng
-------------[2] Phương trình căn √︀f(x) = g(x) và mở rộng
-------------[3] Phương trình căn thức có tham số
-------------[4] Phương trình căn thức (dạng khác)
-------------[5] Phương trình khác quy về phương trình bậc 2
-------------[6] Toán hình, toán thực tế ứng dụng phương trình quy về bậc 2
-------[8] Đại số tổ hợp
----------[1] Quy tắc cộng-quy tắc nhân
-------------[1] Bài toán chỉ sử dụng quy tắc cộng
-------------[2] Bài toán chỉ sử dụng quy tắc nhân
-------------[3] Bài toán kết hợp quy tắc cộng và quy tắc nhân
-------------[4] Bài toán dùng quy tắc bù trừ
-------------[5] Bài toán đếm số tự nhiên
-------------[6] Sơ đồ hình cây
----------[2] Hoán vị. Chỉnh hợp. Tổ hợp
-------------[1] Lý thuyết tổng hợp về P,C,A
-------------[2] Bài toán có biểu thức P,C,A
-------------[3] Bài toán đếm số tự nhiên
-------------[4] Bài toán chọn người
-------------[5] Bài toán chọn đối tượng khác
-------------[6] Bài toán có yếu tố hình học
-------------[7] Bài toán xếp chỗ (không tròn, không lặp)
-------------[8] Hoán vị bàn tròn
-------------[9] Hoán vị lặp
----------[3] Nhị thức Newton (Kể cả chuyên đề)
-------------[1] Các câu hỏi lý thuyết tổng hợp
-------------[2] Khai triển một nhị thức Newton
-------------[3] Tìm hệ số, số hạng trong khai triển bằng tam giác Pascal
-------------[4] Tìm hệ số, số hạng trong khai triển
-------------[5] Tính tổng nhờ khai triển nhị thức Newton
-------------[6] Toán tổ hợp có dùng nhị thức Newton
-------[0] Xác suất
----------[1] Không gian mẫu và biến cố
-------------[1] Các câu hỏi lý thuyết tổng hợp
-------------[2] Mô tả không gian mẫu, biến cố
-------------[3] Đếm phần tử không gian mẫu, biến cố
----------[2] Xác suất của biến cố
-------------[1] Các câu hỏi lý thuyết tổng hợp
-------------[2] Liên quan xúc xắc, đồng tiền (PP liệt kê)
-------------[3] Liên quan việc sắp xếp chỗ
-------------[4] Liên quan việc chọn người
-------------[5] Liên quan việc chọn đối tượng khác
-------------[6] Liên quan hình học
-------------[7] Liên quan việc đếm số
-------------[8] Liên quan bàn tròn hoặc hoán vị lặp 
-------------[9] Liên quan vấn đề khác
----[H] Hình học 10
-------[4] Hệ thức lượng trong tam giác
----------[1] Giá trị lượng giác của góc (0 − 180∘)
-------------[1] Xét dấu của biểu thức lượng giác
-------------[2] Tính các giá trị lượng giác
-------------[3] Biến đổi, rút gọn biểu thức lượng giác
----------[2] Định lý sin và định lý côsin
-------------[1] Bài toán chỉ dùng định lý Sin, Côsin
-------------[2] Bài toán có dùng công thức diện tích
-------------[3] Biến đổi, rút gọn biểu thức
-------------[4] Nhận dạng tam giác
----------[3] Giải tam giác và ứng dụng thực tế
-------------[1] Giải tam giác
-------------[2] Các ứng dụng thực tế
-------[5] Véctơ (chưa xét tọa độ)
----------[1] Khái niệm véctơ
-------------[1] Xác định một véctơ
-------------[2] Xét phương và hướng của các véctơ
-------------[3] Hai véctơ bằng nhau
-------------[4] Hai véctơ đối nhau
-------------[5] Độ dài của một véctơ
-------------[6] Toán thực tế áp dụng véctơ
----------[2] Tổng và hiệu của hai véctơ
-------------[1] Tính toán, thu gọn hiệu các véctơ
-------------[2] Tính đúng-sai của 1 đẳng thức véctơ
-------------[3] Tìm điểm nhờ đẳng thức véctơ
-------------[4] Tính độ dài của véctơ tổng, hiệu
-------------[5] Cực trị hình học
-------------[6] Toán thực tế áp dụng tổng hiệu véctơ
----------[3] Tích của một số với véctơ
-------------[1] Xác định k.→−a và độ dài của nó
-------------[2] Biến đổi, thu gọn 1 đẳng thức véctơ
-------------[3] Tìm điểm nhờ đẳng thức véctơ
-------------[4] Sự cùng phương của 2 véctơ và ứng dụng
-------------[5] Phân tích 1 véctơ theo 2 véctơ không cùng phương
-------------[6] Tính độ dài của véctơ tổng, hiệu
-------------[7] Tập hợp điểm
-------------[8] Cực trị hình học
-------------[9] Toán thực tế áp dụng tích 1 số với véctơ
----------[4] Tích vô hướng (chưa xét tọa độ)
-------------[1] Tích vô hướng, góc giữa 2 véctơ
-------------[2] Tìm góc nhờ tích vô hướng
-------------[3] Đẳng thức về tích vô hướng hoặc độ dài
-------------[4] Điều kiện vuông góc
-------------[5] Các bài toán tìm điểm và tập hợp điểm
-------------[6] Cực trị và chứng minh bất đẳng thức
-------------[7] Toán thực tế áp dụng tích vô hướng
-------[9] Phương pháp toạ độ trong mặt phẳng (Oxy)
----------[1] Toạ độ của véctơ
-------------[1] Tọa độ điểm, độ dài đại số của véctơ trên 1 trục
-------------[2] Phép toán véctơ (tổng, hiệu, tích với số) trong Oxy
-------------[3] Tọa độ điểm và véctơ trên hệ trục Oxy
-------------[4] Sự cùng phương của 2 véctơ và ứng dụng
-------------[5] Phân tích một véctơ theo 2 véctơ không cùng phương
-------------[6] Toán thực tế dùng hệ toạ độ
----------[2] Tích vô hướng (theo tọa độ)
-------------[1] Tích vô hướng, góc giữa 2 véctơ
-------------[2] Tìm góc nhờ tích vô hướng
-------------[3] Đẳng thức về tích vô hướng hoặc độ dài
-------------[4] Điều kiện vuông góc
-------------[5] Các bài toán tìm điểm và tập hợp điểm
-------------[6] Cực trị và chứng minh bất đẳng thức
-------------[7] Toán thực tế, liên môn
----------[3] Đường thẳng trong mặt phẳng toạ độ
-------------[1] Điểm, véctơ, hệ số góc của đường thẳng
-------------[2] Phương trình đường thẳng
-------------[3] Vị trí tương đối giữa hai đường thẳng
-------------[4] Bài toán về góc giữa hai đường thẳng
-------------[5] Bài toán về khoảng cách
-------------[6] Bài toán tìm điểm
-------------[7] Bài toán dùng cho tam giác, tứ giác
-------------[8] Bài toán thực tế, PP tọa độ hóa
-------------[9] Bài toán có dùng PT chính tắc
----------[4] Đường tròn trong mặt phẳng toạ độ
-------------[1] Tìm tâm, bán kính và điều kiện là đường tròn
-------------[2] Phương trình đường tròn
-------------[3] Phương trình tiếp tuyến của đường tròn
-------------[4] Vị trí tương đối liên quan đường tròn
-------------[5] Toán tổng hợp đường thẳng và đường tròn
-------------[6] Bài toán dùng cho tam giác, tứ giác
-------------[7] Bài toán thực tế, PP tọa độ hóa
----------[5] Ba đường conic trong mặt phẳng toạ độ và Chuyên đề 3
-------------[1] Xác định các yếu tố của elip
-------------[2] Phương trình chính tắc của elip
-------------[3] Bài toán điểm trên elip
-------------[4] Xác định các yếu tố của hypebol
-------------[5] Phương trình chính tắc của hypebol
-------------[6] Bài toán điểm trên hypebol
-------------[7] Xác định các yếu tố của parabol
-------------[8] Phương trình chính tắc của parabol
-------------[9] Bài toán điểm trên parabol
-------------[0] Bài toán tổng hợp/thực tế, PP tọa độ hóa 3 đường conic
----[C] Chuyên đề 10
-------[1] Hệ phương trình bậc nhất 3 ẩn và ứng dụng
----------[1] Hệ phương trình bậc nhất 3 ẩn và ứng dụng
-------------[1] Các khái niệm về Hệ PT bậc nhất 3 ẩn
-------------[2] Giải Hệ PT bậc nhất 3 ẩn
-------------[3] Toán thực tế ứng dụng Hệ PT bậc nhất 3 ẩn
-------[2] Phương pháp quy nạp toán học
----------[1] Phương pháp quy nạp toán học
-------------[1] Quy nạp chứng minh các đẳng thức/công thức/chia hết
-------------[2] Quy nạp chứng minh các bất đẳng thức
-[1] Lớp 11
----[D] Đại số & Giải tích/Thống kê 11
-------[1] Hàm số lượng giác và phương trình lượng giác
----------[1] Góc lượng giác
-------------[1] Câu hỏi lý thuyết
-------------[2] Chuyển đổi đơn vị độ và radian
-------------[3] Số đo của một góc lượng giác
-------------[4] Độ dài của một cung tròn
-------------[5] Đường tròn lượng giác và ứng dụng
-------------[6] Toán thực tế áp dụng góc lượng giác
----------[2] Giá trị lượng giác của một góc lượng giác
-------------[1] Câu hỏi lý thuyết
-------------[2] Xét dấu giá trị lượng giác. Tính giá trị lượng giác của một góc
-------------[3] Biến đổi, rút gọn biểu thức lượng giác; chứng minh đẳng thức lượng giác
-------------[4] Các góc lượng giác có liên quan đặc biệt: bù nhau, phụ nhau, đối nhau, hơn kém nhau π
-------------[5] Toán thực tế áp dụng giá trị của một góc lượng giác
----------[3] Các công thức lượng giác
-------------[1] Câu hỏi lý thuyết
-------------[2] Áp dụng công thức cộng
-------------[3] Áp dụng công thức nhân đôi - hạ bậc
-------------[4] Áp dụng công thức biến đổi tích <-> tổng
-------------[5] Kết hợp nhiều công thức lượng giác
-------------[6] Nhận dạng tam giác
-------------[7] Toán thực tế áp dụng công thức lượng giác
----------[4] Hàm số lượng giác và đồ thị
-------------[1] Câu hỏi lý thuyết
-------------[2] Tìm tập xác định
-------------[3] Xét tính đơn điệu
-------------[4] Xét tính chẵn, lẻ
-------------[5] Xét tính tuần hoàn, tìm chu kỳ
-------------[6] Tìm tập giá trị và min, max
-------------[7] Bảng biến thiên và đồ thị
-------------[8] Toán thực tế áp dụng hàm số lượng giác
----------[5] Phương trình lượng giác cơ bản
-------------[1] Câu hỏi lý thuyết. Khái niệm phương trình tương đương
-------------[2] Điều kiện có nghiệm
-------------[3] Phương trình cơ bản dùng Radian
-------------[4] Phương trình cơ bản dùng Độ
-------------[5] Phương trình đưa về dạng cơ bản
-------------[6] Toán thực tế áp dụng phương trình lượng giác
----------[6] [Giảm] Phương trình lượng giác thường gặp
-------------[1] Phương trình bậc n theo một hàm số lượng giác
-------------[2] Phương trình đẳng cấp bậc n đối với sinx và cosx
-------------[3] Phương trình bậc nhất đối với sinx và cosx
-------------[4] Phương trình đối xứng, phản đối xứng
-------------[5] Phương trình lượng giác không mẫu mực
-------------[6] Phương trình lượng giác có chứa ẩn ở mẫu số
-------------[7] Phương trình lượng giác có chứa tham số
-------------[8] Toán thực tế áp dụng phương trình lượng giác thường gặp
-------[2] Dãy số. Cấp số cộng. Cấp số nhân
----------[1] Dãy số
-------------[1] Câu hỏi lý thuyết
-------------[2] Số hạng tổng quát, biểu diễn dãy số
-------------[3] Tìm số hạng cụ thể của dãy số
-------------[4] Dãy số tăng, dãy số giảm
-------------[5] Dãy số bị chặn
-------------[6] Toán thực tế áp dụng dãy số
----------[2] Cấp số cộng
-------------[1] Câu hỏi lý thuyết
-------------[2] Nhận diện cấp số cộng, công sai d
-------------[3] Số hạng tổng quát của cấp số cộng
-------------[4] Tìm số hạng cụ thể trong cấp số cộng
-------------[5] Điều kiện để dãy số là cấp số cộng
-------------[6] Tính tổng của cấp số cộng
-------------[7] Toán thực tế áp dụng cấp số cộng
----------[3] Cấp số nhân
-------------[1] Câu hỏi lý thuyết
-------------[2] Nhận diện cấp số nhân, công bội q
-------------[3] Số hạng tổng quát của cấp số nhân
-------------[4] Tìm số hạng cụ thể trong cấp số nhân
-------------[5] Điều kiện để dãy số là cấp số nhân
-------------[6] Tính tổng của cấp số nhân
-------------[7] Kết hợp cấp số nhân và cấp số cộng
-------------[8] Toán thực tế áp dụng cấp số nhân
-------[3] Giới hạn. Hàm số liên tục
----------[1] Giới hạn của dãy số
-------------[1] Câu hỏi lý thuyết
-------------[2] Phương pháp đặt thừa số chung (lim hữu hạn)
-------------[3] Phương pháp lượng liên hợp (lim hữu hạn)
-------------[4] Giới hạn vô cực
-------------[5] Cấp số nhân lùi vô hạn
-------------[6] Toán thực tế áp dụng giới hạn của dãy số
----------[2] Giới hạn của hàm số
-------------[1] Câu hỏi lý thuyết
-------------[2] Thay số trực tiếp
-------------[3] PP đặt thừa số chung, kết quả hữu hạn
-------------[4] PP đặt thừa số chung, kết quả vô cực
-------------[5] PP lượng liên hợp, kết quả hữu hạn
-------------[6] PP lượng liên hợp, kết quả vô cực
-------------[7] Giới hạn một bên
-------------[8] Toán thực tế áp dụng giới hạn của hàm số
----------[3] Hàm số liên tục
-------------[1] Câu hỏi lý thuyết
-------------[2] Tính liên tục thể hiện qua đồ thị
-------------[3] Hàm số liên tục tại một điểm
-------------[4] Hàm số liên tục trên khoảng, đoạn
-------------[5] Bài toán phương trình có nghiệm
-------------[6] Toán thực tế áp dụng hàm số liên tục
-------[5] Các số đặc trưng đo xu thế trung tâm cho mẫu số liệu ghép nhóm
----------[1] Số trung bình và mốt của mẫu số liệu ghép nhóm
-------------[1] Câu hỏi lý thuyết
-------------[2] Mẫu số liệu ghép nhóm
-------------[3] Số trung bình
-------------[4] Mốt
----------[2] Trung vị và tứ phân vị của mẫu số liệu ghép nhóm
-------------[1] Câu hỏi lý thuyết
-------------[2] Trung vị
-------------[3] Tứ phân vị
-------[6] Hàm số mũ và hàm số lôgarít
----------[1] Phép tính luỹ thừa
-------------[1] Tính giá trị của biểu thức chứa lũy thừa
-------------[2] Biến đổi, rút gọn biểu thức chứa lũy thừa
-------------[3] Điều kiện cho luỹ thừa, căn thức
-------------[4] So sánh các lũy thừa
----------[2] Phép tính lôgarít
-------------[1] Tính giá trị biểu thức chứa lôgarít
-------------[2] Biến đổi, biểu diễn biểu thức chứa lôgarít
-------------[3] Rút gọn, chứng minh biểu thức lôgarít
-------------[4] Số e và bài toán lãi kép
-------------[5] Toán thực tế áp dụng phép tính lôgarít
----------[3] Hàm số mũ. Hàm số lôgarít
-------------[1] Câu hỏi lý thuyết hàm số lũy thừa, mũ, lôgarít
-------------[2] Tập xác định của hàm số
-------------[3] Sự biến thiên và đồ thị của hàm số mũ, lôgarít
-------------[4] So sánh các luỹ thừa và lôgarít
-------------[5] Toán thực tế áp dụng hàm số mũ, lôgarít
----------[4] Phương trình, bất phương trình mũ và lôgarít
-------------[1] Điều kiện có nghiệm
-------------[2] Phương trình mũ, lôgarít cơ bản
-------------[3] Bất phương trình mũ, lôgarít cơ bản
-------------[4] Phương trình mũ, lôgarít đưa về cùng cơ số
-------------[5] Bất phương trình mũ, lôgarít đưa về cùng cơ số
-------------[6] Toán thực tế áp dụng phương trình mũ, lôgarít
----------[5] [Giảm] Các phương pháp giải được giảm tải
-------------[1] Phương pháp đặt ẩn phụ cho PT mũ, lôgarít
-------------[2] Phương pháp lôgarít hóa, mũ cho PT mũ, lôgarít
-------------[3] Phương pháp hàm số, đánh giá cho PT mũ, lôgarít
-------------[4] Hệ PT mũ, lôgarít
-------------[5] Toán thực tế áp dụng phương trình mũ, lôgarít
-------[7] Đạo hàm
----------[1] Đạo hàm
-------------[1] Tính đạo hàm bằng định nghĩa
-------------[2] Số gia hàm số, số gia biến số
-------------[3] Ý nghĩa Hình học của đạo hàm
-------------[4] Ý nghĩa Vật lý của đạo hàm
-------------[5] Toán thực tế khác áp dụng định nghĩa đạo hàm
----------[2] Các quy tắc đạo hàm
-------------[1] Tính đạo hàm
-------------[2] Đẳng thức có y và y′
-------------[3] Tiếp tuyến tại một điểm
-------------[4] Tiếp tuyến biết trước hệ số góc
-------------[5] Tiếp tuyến chưa biết tiếp điểm và hệ số góc
-------------[6] Giới hạn hàm số lượng giác, hàm số mũ, lôgarít
-------------[7] Dùng đạo hàm cho nhị thức Newton
-------------[8] Toán thực tế áp dụng quy tắc đạo hàm
----------[3] Đạo hàm cấp hai
-------------[1] Tính đạo hàm cấp hai
-------------[2] Đẳng thức có y và (y′, y′′)
-------------[3] Toán thực tế và Ý nghĩa Vật lý của đạo hàm cấp hai
-------[9] Xác suất
----------[1] Biến cố giao và quy tắc nhân xác suất
-------------[1] Câu hỏi lí thuyết
-------------[2] Xác định và đếm số phần tử biến cố giao
-------------[3] Công thức nhân xác suất cho 2 biến cố độc lập
-------------[4] Tính xác suất biến cố giao bằng sơ đồ hình cây
----------[2] Biến cố hợp và quy tắc cộng xác suất
-------------[1] Câu hỏi lí thuyết
-------------[2] Mô tả biến cố hợp
-------------[3] Tính và xác định số phần tử biến cố hợp
-------------[4] Quy tắc cộng cho hai biến cố xung khắc
-------------[5] Quy tắc cộng cho hai biến cố bất kì
-------------[6] Tính xác suất bằng sơ đồ hình cây
----[H] Hình học 11
-------[4] Đường thẳng, mặt phẳng. Quan hệ song song trong không gian
----------[1] Điểm, đường thẳng và mặt phẳng
-------------[1] Câu hỏi lý thuyết
-------------[2] Hình biểu diễn của một hình không gian
-------------[3] Tìm giao tuyến của hai mặt phẳng
-------------[4] Tìm giao điểm của đường thẳng và mặt phẳng
-------------[5] Xác định thiết diện
-------------[6] Ba điểm thẳng hàng, ba đường thẳng đồng quy
-------------[7] Toán thực tế áp dụng điểm, đường thẳng và mặt phẳng
----------[2] Hai đường thẳng song song
-------------[1] Câu hỏi lý thuyết
-------------[2] Hai đường thẳng song song
-------------[3] Tìm giao tuyến bằng cách kẻ song song
-------------[4] Tìm giao điểm của đường thẳng và mặt phẳng
-------------[5] Xác định thiết diện bằng cách kẻ song song
-------------[6] Ba điểm thẳng hàng
-------------[7] Bài toán quỹ tích và điểm cố định
-------------[8] Toán thực tế áp dụng hai đường thẳng song song
----------[3] Đường thẳng và mặt phẳng song song
-------------[1] Câu hỏi lý thuyết
-------------[2] Đường thẳng song song với mặt phẳng
-------------[3] Tìm giao tuyến bằng cách kẻ song song
-------------[4] Tìm giao điểm của đường thẳng và mặt phẳng
-------------[5] Xác định thiết diện bằng cách kẻ song song
-------------[6] Ba điểm thẳng hàng
-------------[7] Bài toán quỹ tích và điểm cố định
-------------[8] Toán thực tế áp dụng đường thẳng song song mặt phẳng
----------[4] Hai mặt phẳng song song
-------------[1] Câu hỏi lý thuyết
-------------[2] Hai mặt phẳng song song
-------------[3] Chứng minh đường thẳng song song mặt phẳng
-------------[4] Xác định mặt phẳng đi qua một điểm và song song với một mặt phẳng
-------------[5] Xác định mặt phẳng chứa đường thẳng (hoặc đi qua hai điểm) và song song với một mặt phẳng
-------------[6] Bài toán tổng hợp
-------------[7] Toán thực tế áp dụng hai mặt phẳng song song
----------[5] Hình lăng trụ và hình hộp (xiên)
-------------[1] Câu hỏi lý thuyết
-------------[2] Bài toán về hình lăng trụ (xiên)
-------------[3] Bài toán về hình hộp (xiên)
-------------[4] Toán thực tế áp dụng hình lăng trụ và hình hộp
----------[6] Phép chiếu song song
-------------[1] Câu hỏi lý thuyết
-------------[2] Hình biểu diễn của một hình không gian
-------------[3] Xác định yếu tố song song
-------------[4] Xác định phương chiếu
-------------[5] Tính tỉ số đoạn thẳng, diện tích qua phép chiếu
-------[8] Quan hệ vuông góc trong không gian
----------[1] Hai đường thẳng vuông góc
-------------[1] Câu hỏi lí thuyết
-------------[2] Xác định hai đường thẳng vuông góc
-------------[3] Tìm góc giữa hai đường thẳng
-------------[4] Toán thực tế áp dụng hai đường thẳng vuông góc
----------[2] Đường thẳng vuông góc với mặt phẳng
-------------[1] Câu hỏi lí thuyết
-------------[2] Xác định hoặc chứng minh đường thẳng và mặt phẳng vuông góc
-------------[3] Xác định hoặc chứng minh hai đường thẳng vuông góc
-------------[4] Dựng mặt phẳng, tìm thiết diện
-------------[5] Hình chiếu vuông góc của một hình trên mặt phẳng (tìm điểm, tìm đoạn thẳng, tính diện tích)
-------------[6] Toán thực tế áp dụng đường thẳng vuông góc mặt phẳng
----------[3] Phép chiếu vuông góc
-------------[1] Lý thuyết về phép chiếu vuông góc
-------------[2] Hình chiếu vuông góc của đa giác trên mặt phẳng
-------------[3] Các bài toán thực tế áp dụng phép chiếu vuông góc
----------[4] Hai mặt phẳng vuông góc
-------------[1] Câu hỏi lí thuyết
-------------[2] Xác định/chứng minh đường thẳng vuông góc mặt phẳng, mặt phẳng vuông góc
-------------[3] Xác định góc giữa hai mặt phẳng
-------------[4] Dựng mặt phẳng, thiết diện
-------------[5] Nhận dạng và tính toán liên quan các hình thông dụng
-------------[6] Bài toán cho trước góc giữa d và (P)
-------------[7] Toán thực tế áp dụng hai mặt phẳng vuông góc
----------[5] Khoảng cách
-------------[1] Câu hỏi lí thuyết
-------------[2] Khoảng cách giữa 2 điểm, từ 1 điểm đến 1 đường thẳng
-------------[3] Khoảng cách từ một điểm đến một mặt phẳng
-------------[4] Khoảng cách giữa hai đường thẳng chéo nhau
-------------[5] Đường vuông góc chung của hai đường thẳng chéo nhau
-------------[6] Toán thực tế áp dụng khoảng cách
----------[6] Góc giữa đường thẳng và mặt phẳng. Góc nhị diện
-------------[1] Góc giữa đường thẳng và mặt phẳng
-------------[2] Góc nhị diện, góc phẳng nhị diện
-------------[3] Góc giữa 2 mặt phẳng, biết trước góc (d,(P))
-------------[4] Khoảng cách giữa điểm, đường, biết trước góc (d,(P))
-------------[5] Khoảng cách giữa điểm - mặt phẳng, biết trước góc (d,(P))
-------------[6] Khoảng cách giữa 2 đường chéo nhau, biết trước góc (d,(P))
-------------[7] Toán thực tế về góc đường thẳng, mặt phẳng, góc nhị diện
----------[7] Hình lăng trụ đứng. Hình chóp đều. Thể tích
-------------[1] Câu hỏi lí thuyết, công thức
-------------[2] Thể tích khối chóp tam giác
-------------[3] Thể tích khối chóp tứ giác
-------------[4] Thể tích khối lăng trụ tam giác
-------------[5] Thể tích khối lăng trụ tứ giác
-------------[6] Thể tích khối chóp cụt và các khối khác
-------------[7] Tỉ số thể tích
-------------[8] Ứng dụng thể tích tính góc, khoảng cách
-------------[9] Toán thực tế hình lăng trụ đứng, chóp đều, thể tích
----[C] Chuyên đề 11
-------[1] Phép biến hình phẳng
----------[1] Phép biến hình, phép dời hình và hai hình bằng nhau
-------------[1] Câu hỏi lý thuyết
-------------[2] Bài toán xác định một phép đặt tương ứng có là phép dời hình hay không?
-------------[3] Xác định ảnh khi thực hiện phép dời hình
----------[2] Phép tịnh tiến
-------------[1] Câu hỏi lý thuyết
-------------[2] Tìm ảnh hoặc tạo ảnh khi thực hiện phép tịnh tiến
-------------[3] Ứng dụng phép tịnh tiến
----------[3] Phép đối xứng trục
-------------[1] Câu hỏi lý thuyết
-------------[2] Tìm ảnh hoặc tạo ảnh khi thực hiện phép đối xứng trục
-------------[3] Xác định trục đối xứng và số trục đối xứng của một hình
-------------[4] Ứng dụng phép đối xứng trục
----------[4] Phép đối xứng tâm
-------------[1] Câu hỏi lý thuyết
-------------[2] Tìm ảnh, tạo ảnh khi thực hiện phép đối xứng tâm
-------------[3] Xác định hình có tâm đối xứng
-------------[4] Ứng dụng phép đối xứng tâm
----------[5] Phép quay
-------------[1] Câu hỏi lý thuyết
-------------[2] Xác định vị trí ảnh của điểm, hình khi thực hiện phép quay cho trước
-------------[3] Tìm tọa độ ảnh của điểm, phương trình của một đường thẳng khi thực hiện phép quay
-------------[4] Ứng dụng phép quay
----------[6] Phép vị tự
-------------[1] Câu hỏi lý thuyết
-------------[2] Xác định ảnh, tạo ảnh khi thực hiện phép vị tự
-------------[3] Tìm tâm vị tự của hai đường tròn
-------------[4] Ứng dụng phép vị tự
----------[7] Phép đồng dạng
-------------[1] Câu hỏi lý thuyết
-------------[2] Xác định ảnh, tạo ảnh khi thực hiện phép đồng dạng
-------[2] Lý thuyết đồ thị
----------[1] Đồ thị
-------------[1] Câu hỏi về đỉnh, cạnh của đồ thị
-------------[2] Câu hỏi về bậc của đồ thị
-------------[3] Câu hỏi tổng hợp
----------[2] Đường đi Euler và Harmilton
-------------[1] Đường đi Euler
-------------[2] Đường đi Harmilton
-------------[3] Câu hỏi tổng hợp
----------[3] Bài toán tìm đường đi ngắn nhất
-------------[1] Bài toán tìm đường đi ngắn nhất
-------------[2] Tổng hợp
-------[3] Một số yếu tố vẽ kỹ thuật
----------[1] Hình biểu diễn của một hình, khối
-------------[1] Lý thuyết về phép chiếu và hình biểu diễn song song
-------------[2] Lý thuyết về phép chiếu vuông góc
-------------[3] Lý thuyết về phép chiếu trục đo
-------------[4] Tổng hợp
----------[2] Bản vẽ kỹ thuật
-------------[1] Lý thuyết cơ bản về bản vẽ kỹ thuật
-------------[2] Phương pháp biểu diễn bản vẽ kỹ thuật
-------------[3] Tổng hợp
-[2] Lớp 12
----[D] Giải tích
-------[1] Ứng dụng đạo hàm để khảo sát hàm số
----------[1] Sự đồng biến và nghịch biến của hàm số
-------------[1] Xét tính đơn điệu của hàm số cho bởi công thức
-------------[2] Xét tính đơn điệu dựa vào bảng biến thiên, đồ thị
-------------[3] Tìm tham số m để hàm số đơn điệu
-------------[4] Ứng dụng tính đơn điệu để chứng minh bất đẳng thức, giải phương trình, bất phương trình, hệ phương trình
-------------[5] Toán thực tế ứng dụng sự đồng biến nghịch biến
----------[2] Cực trị của hàm số
-------------[1] Tìm cực trị của hàm số cho bởi công thức
-------------[2] Tìm cực trị dựa vào BBT, đồ thị
-------------[3] Tìm m để hàm số đạt cực trị tại 1 điểm x0 cho trước
-------------[4] Tìm m để hàm số, đồ thị hàm số bậc ba có cực trị thỏa mãn điều kiện
-------------[5] Tìm m để hàm số, đồ thị hàm số trùng phương có cực trị thỏa mãn điều kiện
-------------[6] Tìm m để hàm số, đồ thị hàm số các hàm số khác có cực trị thỏa mãn điều kiện
-------------[7] Toán thực tế ứng dụng cực trị của hàm số
----------[3] Giá trị lớn nhất và giá trị nhỏ nhất của hàm số
-------------[1] GTLN, GTNN trên đoạn [a;b]
-------------[2] GTLN, GTNN trên khoảng
-------------[3] Sử dụng các đánh giá, bất đẳng thức cổ điển
-------------[4] Ứng dụng GTNN, GTLN trong bài toán phương trình, bất phương trình, hệ phương trình
-------------[5] GTLN, GTNN hàm nhiều biến
-------------[6] Toán thực tế ứng dụng GTLN, GTNN của hàm số
----------[4] Đường tiệm cận
-------------[1] Bài toán xác định các đường tiệm cận của hàm số (không chứa tham số) hoặc biết BBT, đồ thị
-------------[2] Bài toán xác định các đường tiệm cận của hàm số có chứa tham số
-------------[3] Bài toán liên quan đến đồ thị hàm số và các đường tiệm cận
-------------[4] Toán thực tế ứng dụng tiệm cận
----------[5] Khảo sát sự biến thiên và vẽ đồ thị hàm số
-------------[1] Nhận dạng đồ thị
-------------[2] Các phép biến đổi đồ thị
-------------[3] Biện luận số giao điểm dựa vào đồ thị, bảng biến thiên
-------------[4] Sự tương giao của hai đồ thị (liên quan đến tọa độ giao điểm)
-------------[5] Đồ thị của hàm đạo hàm
-------------[6] Phương trình tiếp tuyến của đồ thị hàm số
-------------[7] Điểm đặc biệt của đồ thị hàm số
-------------[8] Toán thực tế ứng dụng khảo sát hàm số
-------[3] Các số đặc trưng đo mức độ phân tán cho mẫu số liệu ghép nhóm
----------[1] Khoảng biến thiên, khoảng tứ phân vị của mẫu số liệu ghép nhóm
-------------[1] Công thức lý thuyết
-------------[2] Tìm khoảng biến thiên
-------------[3] Tìm khoảng tứ phân vị
-------------[4] Câu hỏi tổng hợp
----------[2] Phương sai, độ lệch chuẩn của mẫu số liệu ghép nhóm
-------------[1] Công thức lý thuyết
-------------[2] Tìm phương sai, độ lệch chuẩn
-------------[3] Câu hỏi tổng hợp
-------[4] Nguyên hàm, tích phân và ứng dụng
----------[1] Nguyên hàm
-------------[1] Công thức lý thuyết
-------------[2] Nguyên hàm cơ bản đa thức, phân thức
-------------[3] Nguyên hàm cơ bản hàm lượng giác
-------------[4] Nguyên hàm cơ bản hàm mũ, luỹ thừa
-------------[5] Phương pháp đổi biến số cơ bản
-------------[6] Toán thực tế áp dụng nguyên hàm
----------[2] Tích phân
-------------[1] Công thức lý thuyết
-------------[2] Tích phân cơ bản đa thức, phân thức
-------------[3] Tích phân cơ bản hàm lượng giác
-------------[4] Tích phân cơ bản hàm mũ, luỹ thừa
-------------[5] Phương pháp đổi biến số cơ bản
-------------[6] Toán thực tế áp dụng nguyên hàm
----------[3] Ứng dụng thực tế và hình học của tích phân
-------------[1] Diện tích hình phẳng được giới hạn bởi các đồ thị
-------------[2] Bài toán thực tế sử dụng diện tích hình phẳng
-------------[3] Thể tích giới hạn bởi các đồ thị (tròn xoay)
-------------[4] Thể tích tính theo mặt cắt S(x)
-------------[5] Bài toán thực tế và ứng dụng thể tích tròn xoay, S(x)
-------[6] Một số yếu tố xác suất
----------[1] Xác suất có điều kiện
-------------[1] Công thức lý thuyết
-------------[2] Tính xác suất có điều kiện bằng công thức
-------------[3] Tính xác suất có điều kiện bằng sơ đồ cây
-------------[4] Bài toán tổng hợp
----------[2] Công thức xác suất toàn phần. Công thức Bayes
-------------[1] Công thức lý thuyết
-------------[2] Tính xác suất bằng công thức xác suất toàn phần
-------------[3] Tính xác suất bằng công thức xác suất Bayes
-------------[4] Bài toán tổng hợp
----[H] Hình học
-------[2] Tọa độ của véc-tơ trong không gian
----------[1] Véc-tơ và các phép toán véc-tơ trong không gian (chưa toạ độ hoá)
-------------[1] Công thức lý thuyết
-------------[2] Tổng, hiệu, tích một số với véc-tơ
-------------[3] Tích vô hướng và ứng dụng
-------------[4] Toán thực tế áp dụng các phép toán véc-tơ
----------[2] Toạ độ của véc-tơ và các công thức
-------------[1] Công thức lý thuyết
-------------[2] Tìm tọa độ điểm
-------------[3] Tìm tọa độ véc-tơ
-------------[4] Công thức toạ độ của tích vô hướng và ứng dụng
-------------[5] Công thức toạ độ của tích có hướng và ứng dụng
-------------[6] Toán thực tế áp dụng các phép toán toạ độ hoá véc-tơ}
-------[5] Phương trình mặt phẳng, đường thẳng, mặt cầu trong không gian Oxyz
----------[1] Phương trình mặt phẳng
-------------[1] Câu hỏi lý thuyết
-------------[2] Xác định véc-tơ pháp tuyến, cặp véc-tơ chỉ phương
-------------[3] Viết phương trình tổng quát mặt phẳng
-------------[4] Vị trí tương đối giữa hai mặt phẳng (song song, vuông góc)
-------------[5] Khoảng cách điểm tới mặt phẳng
-------------[6] Góc giữa hai mặt phẳng
-------------[7] Toán thực tế áp dụng phương trình mặt phẳng
----------[2] Phương trình đường thẳng trong không gian
-------------[1] Câu hỏi lý thuyết
-------------[2] Xác định véc-tơ chỉ phương, cặp véc-tơ pháp tuyến
-------------[3] Viết phương trình tổng quát, chính tắc, tham số đường thẳng
-------------[4] Vị trí tương đối giữa hai đường thẳng
-------------[5] Vị trí tương đối giữa đường thẳng và mặt phẳng
-------------[6] Khoảng cách điểm tới đường thẳng
-------------[7] Góc giữa hai đường thẳng, đường thẳng và mặt phẳng
-------------[8] Toán thực tế áp dụng phương trình đường thẳng
----------[3] Phương trình mặt cầu trong không gian
-------------[1] Câu hỏi lý thuyết
-------------[2] Xác định tâm, bán kính, đường kính mặt cầu
-------------[3] Viết phương trình tổng quát mặt cầu
-------------[4] Toán thực tế áp dụng phương trình mặt cầu
----[C] Chuyên đề
-------[1] Khối đa diện
----------[1] Khái niệm về khối đa diện
-------------[1] Nhận diện hình đa diện, khối đa diện
-------------[2] Xác định số đỉnh, cạnh, mặt bên của một khối đa diện
-------------[3] Phân chia, lắp ghép các khối đa diện
-------------[4] Phép biến hình trong không gian
----------[2] Khối đa diện lồi và khối đa diện đều
-------------[1] Nhận diện đa diện lồi
-------------[2] Nhận diện loại đa diện đều
-------------[3] Tính chất đối xứng
----------[3] Khái niệm về thể tích của khối đa diện
-------------[1] Diện tích xung quanh, diện tích toàn phần của khối đa diện
-------------[2] Tính thể tích các khối đa diện
-------------[3] Tỉ số thể tích
-------------[4] Các bài toán khác(góc, khoảng cách,...) liên quan đến thể tích khối đa diện
-------------[5] Bài toán thực tế về khối đa diện
-------------[6] Bài toán cực trị
-------[2] Mặt nón, mặt trụ, mặt cầu
----------[1] Khái niệm về mặt tròn xoay
-------------[1] Thể tích khối nón, khối trụ
-------------[2] Diện tích xung quanh, diện tích toàn phần, độ dài đường sinh, chiều cao, bán kính đáy, thiết diện
-------------[3] Khối tròn xoay nội tiếp, ngoại tiếp khối đa diện
-------------[4] Bài toán thực tế về khối nón, khối trụ
-------------[5] Bài toán cực trị về khối nón, khối trụ
----------[2] Mặt cầu
-------------[1] Bài toán sử dụng định nghĩa, tính chất, vị trí tương đối
-------------[2] Khối cầu ngoại tiếp khối đa diện
-------------[3] Khối cầu nội tiếp khối đa diện
-------------[4] Bài toán thực tế về khối cầu
-------------[5] Bài toán cực trị về khối cầu
-------------[6] Bài toán tổng hợp về khối nón, khối trụ, khối cầu
-------[3] Phương pháp tọa độ trong không gian
----------[1] Hệ tọa độ trong không gian
-------------[1] Tìm tọa độ điểm, véc-tơ liên quan đến hệ trục Oxyz
-------------[2] Tích vô hướng và ứng dụng
-------------[3] Phương trình mặt cầu (xác định tâm, bán kính, viết PT mặt cầu đơn giản, vị trí tương đối hai mặt cầu, điểm đến mặt cầu, đơn giản)
-------------[4] Các bài toán cực trị
----------[2] Phương trình mặt phẳng
-------------[1] Tích có hướng và ứng dụng
-------------[2] Xác định VTPT
-------------[3] Viết phương trình mặt phẳng
-------------[4] Tìm tọa độ điểm liên quan đến mặt phẳng
-------------[5] Góc
-------------[6] Khoảng cách
-------------[7] Vị trí tương đối giữa hai mặt phẳng, giữa mặt cầu và mặt phẳng
-------------[8] Các bài toán cực trị
----------[3] Phương trình đường thẳng trong không gian
-------------[1] Xác định VTCP
-------------[2] Viết phương trình đường thẳng
-------------[3] Tìm tọa độ điểm liên quan đến đường thẳng
-------------[4] Góc
-------------[5] Khoảng cách
-------------[6] Vị trí tương đối giữa hai đường thẳng, giữa đường thẳng và mặt phẳng
-------------[7] Bài toán liên quan giữa đường thẳng - mặt phẳng - mặt cầu
-------------[8] Các bài toán cực trị
----------[4] Ứng dụng của phương pháp tọa độ
-------------[1] Bài toán HHKG
-------------[2] Bài toán đại số
-[6] Lớp 6
----[D] Phần số
-------[1] SỐ TỰ NHIÊN
----------[1] TẬP HỢP. PHẦN TỬ CỦA TẬP HỢP
-------------[1] Làm quen với tập hợp
-------------[2] Các kí hiệu
-------------[3] Cách cho tập hợp
-------------[4] Bài Toán Thực Tế
----------[2] TẬP HỢP SỐ TỰ NHIÊN. GHI SỐ TỰ NHIÊN
-------------[1] Tập hợp N và N*
-------------[2] Thứ tự trong tập hợp số tự nhiên
-------------[3] Ghi số tự nhiên
-------------[4] Bài Toán Thực Tế
----------[3] CÁC PHÉP TÍNH TRONG TẬP HỢP SỐ TỰ NHIÊN
-------------[1] Phép cộng và phép nhân
-------------[2] Tính chất của phép cộng và phép nhân số tự nhiên
-------------[3] Phép trừ và phép chia hết
-------------[4] Bài Toán Thực Tế
----------[4] LUỸ THỪA VỚI SỐ MŨ TỰ NHIÊN
-------------[1] Luỹ thừa
-------------[2] Nhân hai luỹ thừa cùng cơ số
-------------[3] Chia hai luỹ thừa cùng cơ số
-------------[4] Bài Toán Thực Tế
----------[5] THỨ TỰ THỰC HIỆN CÁC PHÉP TÍNH
-------------[1] Thứ tự thực hiện các phép tính
-------------[2] Sử dụng máy tính cầm tay
-------------[3] Bài Toán Thực Tế
----------[6] CHIA HẾT VÀ CHIA CÓ DƯ. TÍNH CHẤT CHIA HẾT CỦA MỘT TỔNG
-------------[1] Chia hết và chia có dư
-------------[2] Tính chất chia hết của một tổng
-------------[3] Bài toán thực tế
----------[7] DẤU HIỆU CHIA HẾT CHO 2, CHO 5
-------------[1] Dấu hiệu chia hết cho 2
-------------[2] Dấu hiệu chia hết cho 5
-------------[3] Bài Toán Thực Tế
----------[8] DẤU HIỆU CHIA HẾT CHO 3, CHO 9
-------------[1] Dấu hiệu chia hết cho 9
-------------[2] Dấu hiệu chia hết cho 3
-------------[3] Bài Toán Thực Tế
----------[9] ƯỚC VÀ BỘI
-------------[1] Ước và bội
-------------[2] Cách tìm ước
-------------[3] Cách tìm bội
-------------[4] Bài Toán Thực Tế
----------[0] SỐ NGUYÊN TỐ. HỢP SỐ. PHÂN TÍCH MỘT SỐ RA THỪA SỐ NGUYÊN TỐ
-------------[1] Số nguyên tố. Hợp số
-------------[2] Phân tích một số ra thừa số nguyên tố
-------------[3] Bài Toán thực tế
----------[A] ƯỚC CHUNG. ƯỚC CHUNG LỚN NHẤT
-------------[1] Ước chung
-------------[2] Ước chung lớn nhất
-------------[3] Tìm ước chung lớn nhất bằng cách phân tích các số ra thừa số nguyên tố
-------------[4] Ứng dụng trong rút gọn phân số
-------------[5] Bài Toán Thực Tế
----------[B] BỘI CHUNG. BỘI CHUNG NHỎ NHẤT
-------------[1] Bội chung
-------------[2] Bội chung lớn nhất
-------------[3] Tìm bội chung lớn nhất bằng cách phân tích các số ra thừa số nguyên tố
-------------[4] Ứng dụng trong quy đồng mẫu các phân số
-------------[5] Bài Toán Thực Tế
-------[2] SỐ NGUYÊN
----------[1] Số nguyên âm và tập hợp các số nguyên
-------------[1] Làm quen với số nguyên âm
-------------[2] Tập hợp số nguyên
-------------[3] Biểu diễn số nguyên tố trên trục số
-------------[4] Số đối của một số nguyên
-------------[5] Bài Toán Thực Tế
----------[2] Thứ tự trong tập hợp số nguyên
-------------[1] So sánh hai số nguyên
-------------[2] Thứ tự trong tập hợp số nguyên
-------------[3] Bài Toán Thực Tế
----------[3] Phép cộng và phép trừ hai số nguyên
-------------[1] Cộng hai số nguyên cùng dấu
-------------[2] Cộng hai số nguyên khác dấu
-------------[3] Tính chất phép cộng các số nguyên
-------------[4] Phép trừ hai số nguyên
-------------[5] Quy tắc dấu ngoặc
-------------[6] Bài Toán Thực Tế
----------[4] Phép nhân và phép chia hai số nguyên
-------------[1] Nhân hai số nguyên khác dấu
-------------[2] Nhân hai số nguyên cùng dấu
-------------[3] Tính chất phép nhân các số nguyên
-------------[4] Quan hệ chia hết và phép chia hết trong tập hợp số nguyên (Bài toán tìm x)
-------------[5] Bội và ước của một số nguyên
-------------[6] Bài Toán Thực Tế
-------[3] MỘT SỐ YẾU TỐ THỐNG KÊ
----------[1] Thu thập và phân loại dữ liệu
-------------[1] Thu thập dữ liệu
-------------[2] Phân loại dữ liệu
-------------[3] Tính hợp lí của dữ liệu
-------------[4] Bài Toán Thực Tế
----------[2] Biểu diễn dữ liệu trên bảng
-------------[1] Bảng dữ liệu ban đầu
-------------[2] Bảng thống kê
-------------[3] Bài Toán Thực Tế
----------[3] Biểu đồ tranh
-------------[1] Ôn tập và bổ sung kiến thức
-------------[2] Đọc biểu đồ tranh
-------------[3] Vẽ biểu đồ tranh
-------------[4] Bài Toán Thực Tế
----------[4] Biểu đồ cột - biểu đồ cột kép
-------------[1] Ôn tập biểu đồ cột
-------------[2] Đọc biểu đồ cột
-------------[3] Vẽ biểu đồ cột
-------------[4] Giới thiệu biểu đồ kép
-------------[5] Đọc biểu đồ kép
-------------[6] Vẽ biểu đồ kép
-------------[7] Bài toán thực tế
-------[4] PHÂN SỐ
----------[1] Phân số với tử số và mẫu số là số nguyên
-------------[1] Mở rộng khái niệm phân số
-------------[2] Phân số bằng nhau
-------------[3] Biểu diễn số nguyên ở dạng phân số
-------------[4] Bài Toán Thực Tế
----------[2] Tính chất cơ bản của phân số
-------------[1] Tính chất 1 (nhân tử mẫu cùng một số nguyên)
-------------[2] Tính chất 2 (chia tử mẫu cùng một số nguyên)
-------------[3] Bài Toán Thực Tế
----------[3] So sánh phân số
-------------[1] So sánh hai phân số cùng mẫu
-------------[2] So sánh hai phân số khác mẫu
-------------[3] Áp dụng quy tắc so sánh phân số
-------------[4] Bài Toán Thực Tế
----------[4] Phép cộng và phép trừ phân số
-------------[1] Phép cộng hai phân số
-------------[2] Một số tính chất của phép cộng phân số
-------------[3] Số đối
-------------[4] Phép trừ hai phân số
-------------[5] Bài Toán Thực Tế
----------[5] Phép nhân và phép chia phân số
-------------[1] Nhân hai phân số
-------------[2] Một số tính chất của phép nhân phân số
-------------[3] Chia phân số
-------------[4] Bài Toán Thực Tế
----------[6] Giá trị phân số của một số
-------------[1] Tính giá trị phân số của một số
-------------[2] Tìm một số khi biết giá trị phân số của số đó
-------------[3] Bài Toán Thực Tế
----------[7] Hỗn số
-------------[1] Hỗn số
-------------[2] Đổi hỗn số ra phân số
-------------[3] Bài Toán Thực Tế
-------[5] SỐ THẬP PHÂN
----------[1] Số thập phân
-------------[1] Số thập phân âm
-------------[2] Số đối của một số thập phân
-------------[3] So sánh hai số thập phân
-------------[4] Bài Toán Thực Tế
----------[2] Các phép tính với số thập phân
-------------[1] Cộng, trừ hai số thập phân
-------------[2] Nhân, chia hai số thập phân dương
-------------[3] Nhân, chia hai số thập phân có dấu bất kì
-------------[4] Tính chất của các phép tính với số thập phân
-------------[5] Bài Toán Thực Tế
----------[3] Làm tròn số thập phân và ước lượng kết quả
-------------[1] Làm tròn số thập phân
-------------[2] Ước lượng kết quả
-------------[3] Bài Toán Thực Tế
----------[4] Tỉ số và tỉ số phần trăm
-------------[1] Tỉ số của hai đại lượng
-------------[2] Tỉ số phần trăm của hai đại lượng
-------------[3] Bài Toán Thực Tế
----------[5] Bài toán về tỉ số phần trăm
-------------[1] Tìm giá trị phần trăm của một số
-------------[2] Tìm một số khi biết giá trị phần trăm của số đó
-------------[3] Sử dụng tỉ số phần trăm trong thực tế
-------[6] MỘT SỐ YẾU TỐ XÁC SUẤT
----------[1] Phép thử nghiệm. Sự kiện
-------------[1] Phép thử nghiệm
-------------[2] Sự kiện
-------------[3] Bài Toán Thực Tế
----------[2] Xác suất thực nghiệm
-------------[1] Khả năng xảy ra của một sự kiện
-------------[2] Xác suất thực nghiệm
-------------[3] Bài Toán Thực Tế
----[H] Hình học
-------[1] CÁC HÌNH PHẲNG TRONG THỰC TIỄN
----------[1] Hình Vuông - Tam Giác Đều - Lục Giác Đều
-------------[1] Hình vuông
-------------[2] Tam giác đều
-------------[3] Lục giác đều
-------------[4] Bài Toán Thực Tế
----------[2] Hình Chữ Nhật - Hình Thoi - Hình Bình Hành - Hình Thang Cân
-------------[1] Hình chữ nhật
-------------[2] Hình thoi
-------------[3] Hình bình hành
-------------[4] Hình thang cân
-------------[5] Bài Toán Thực Tế
----------[3] Chu Vi và Diện Tích của một số hình trong thực tiễn
-------------[1] Chu vi và diện tích hình chữ nhật, hình vuông, hình tam giác, hình thang
-------------[2] Tính chu vi và diện tích hình bình hành, hình thoi
-------------[3] Tính chu vi và diện tích một số hình trong thực tiễn
-------[2] TÍNH ĐỐI XỨNG CỦA HÌNH PHẲNG TRONG THẾ GIỚI TỰ NHIÊN
----------[1] Hình có trục đối xứng
-------------[1] Hình có trục đối xứng. Trục đối xứng
-------------[2] Nhận biết hình phẳng trong tự nhiên có trục đối xứng
-------------[3] Bài Toán Thực Tế
----------[2] Hình có tâm đối xứng
-------------[1] Hình có tâm đối xứng. Tâm đối xứng
-------------[2] Nhận biết hình phẳng trong tự nhiên có tâm đối xứng
-------------[3] Bài Toán Thực Tế
----------[3] Vai trò của tính đối xứng trong thế giới tự nhiên
-------------[1] Vẻ đẹp của thế giới tự nhiên biểu hiện qua tính đối xứng
-------------[2] Tính đối xứng trong khoa học kĩ thuật và đời sống
-------------[3] Bài Toán Thực Tế
-------[3] CÁC HÌNH HÌNH HỌC CƠ BẢN
----------[1] Điểm và đường thẳng
-------------[1] Điểm
-------------[2] Đường thẳng
-------------[3] Vẽ đường thẳng
-------------[4] Điểm thuộc đường thẳng. Điểm không thuộc đường thẳng
-------------[5] Bài Toán Thực Tế
----------[2] Ba điểm thẳng hàng. Ba điểm không thẳng hàng
-------------[1] Ba điểm thẳng hàng
-------------[2] Quan hệ giữa ba điểm thẳng hàng
-------------[3] Bài Toán Thực Tế
----------[3] Hai đường thẳng cắt nhau, song song. Tia
-------------[1] Hai đường thẳng cắt nhau, song song
-------------[2] Tia
-------------[3] Bài Toán Thực Tế
----------[4] Đoạn thẳng. Độ dài đoạn thẳng
-------------[1] Đoạn thẳng
-------------[2] Độ dài đoạn thẳng
-------------[3] So sánh hai đoạn thẳng
-------------[4] Một số dụng cụ đo độ dài
-------------[5] Bài Toán Thực Tế
----------[5] Trung điểm của đoạn thẳng
-------------[1] Trung điểm của đoạn thẳng
-------------[2] Cách vẽ trung điểm của đoạn thẳng
-------------[3] Bài Toán Thực Tế
----------[6] Góc
-------------[1] Góc
-------------[2] Cách vẽ góc
-------------[3] Góc bẹt
-------------[4] Điểm trong góc
-------------[5] Bài Toán Thực Tế
----------[7] Số đo góc. Các góc đặc biệt
-------------[1] Thước đo góc
-------------[2] Cách đo góc. Số đo góc
-------------[3] So sánh hai góc
-------------[4] Các góc đặc biệt
-------------[5] Bài Toán Thực Tế
-[7] Lớp 7
----[D] Phần số
-------[1] SỐ HỮU TỈ
----------[1] TẬP HỢP CÁC SỐ HỮU TỈ
-------------[1] Số Hữu Tỉ
-------------[2] Thứ tự trong tập hợp các Số Hữu Tỉ
-------------[3] Biểu diễn Số Hữu Tỉ trên trục
-------------[4] Số đối của Số Hữu Tỉ
-------------[5] Bài Toán Thực Tế
----------[2] CÁC PHÉP TÍNH VỚI SỐ HỮU TỈ
-------------[1] Cộng và trừ 2 Số Hữu Tỉ
-------------[2] Tính chất của phép cộng 2 Số Hữu Tỉ
-------------[3] Nhân 2 Số Hữu Tỉ
-------------[4] Tính chất của phép nhân 2 Số Hữu Tỉ
-------------[5] Chia 2 Số Hữu Tỉ
-------------[6] Bài Toán Thực Tế
----------[3] LUỸ THỪA MỘT SỐ HỮU TỈ
-------------[1] Luỹ thừa với số mũ tự nhiên
-------------[2] Tích và thương 2 luỹ thừa cùng cơ số
-------------[3] Luỹ thừa của luỹ thừa
-------------[4] Bài toán tìm x (phép nhân và phép chia)
-------------[5] Bài Toán Thực Tế
----------[4] QUY TẮC DẤU NGOẶC VÀ CHUYỂN VẾ
-------------[1] Quy tắc dấu ngoặc
-------------[2] Quy tắc chuyển vế
-------------[3] Thứ tự phép tính
-------------[4] Bài toán tìm x (phép cộng và trừ, kết hợp nhân và chia nếu có)
-------------[5] Bài Toán Thực Tế
-------[2] SỐ THỰC
----------[1] SỐ VÔ TỈ. CĂN BẬC HAI SỐ HỌC
-------------[1] Biểu diễn thập phân của số hữu tỉ
-------------[2] Số vô tỷ
-------------[3] Căn bậc hai số học
-------------[4] Tính căn bậc hai số học bằng máy tính
-------------[5] Bài Toán Thực Tế
----------[2] SỐ THỰC. GIÁ TRỊ TUYỆT ĐỐI CỦA MỘT SỐ THỰC
-------------[1] Số thực và tập hợp các số thực
-------------[2] Thứ tự trong tập hợp các số thực
-------------[3] Trục số thực
-------------[4] Số đối của một số thực
-------------[5] Giá trị tuyệt đối của một số thực
-------------[6] Bài toán thực tế
----------[3] LÀM TRÒN SỐ, ƯỚC LƯỢNG KẾT QUẢ
-------------[1] Làm tròn số
-------------[2] Làm tròn số căn cứ vào độ chính xác cho trước
-------------[3] Ước lượng các phép tính
-------------[4] Bài Toán Thực Tế
-------[3] MỘT SỐ YẾU TỐ THỐNG KÊ
----------[1] Thu thập và phân loại dữ liệu.
-------------[1] Thu thập dữ liệu
-------------[2] Phân loại dữ liệu theo các tiêu chí
-------------[3] Tính hợp lí của dữ liệu
----------[2] BIỂU ĐỒ HÌNH QUẠT
-------------[1] Ôn tập về biểu đồ hình quạt tròn
-------------[2] Biểu diễn dữ liệu vào biểu đồ hình quạt tròn
-------------[3] Phân tích dữ liệu trên biểu đồ hình quạt tròn
----------[3] BIỂU ĐỒ ĐOẠN THẲNG
-------------[1] Giới thiệu biểu đồ đoạn thẳng
-------------[2] Vẽ biểu đồ đoạn thẳng
-------------[3] Đọc và phân tích dữ liệu từ biểu đồ đoạn thẳng
-------[4] CÁC ĐẠI LƯỢNG TỈ LỆ
----------[1] Tỉ Lệ Thức - Dãy tỉ số bằng nhau
-------------[1] Tỉ lệ thức
-------------[2] Dãy tỉ số bằng nhau
-------------[3] Tìm x, y, z bằng dãy tỉ số bằng nhau
-------------[4] Bài Toán Thực Tế
----------[2] Đại Lượng Tỉ Lệ Thuận
-------------[1] Đại Lượng Tỉ Lệ Thuận
-------------[2] Tính chất của Đại Lượng Tỉ Lệ Thuận
-------------[3] Các bài toán về Đại Lượng Tỉ Lệ Thuận
-------------[4] Bài Toán Thực Tế
----------[3] Đại lượng tỉ lệ nghịch
-------------[1] Đại lượng tỉ lệ nghịch
-------------[2] Tính chất của các đại lượng tỉ lệ nghịch
-------------[3] Các bài toán về đại lượng tỉ lệ nghịch
-------------[4] Bài Toán Thực Tế
-------[5] BIỂU THỨC ĐẠI SỐ
----------[1] Biểu thức số, Biểu thức đại số
-------------[1] Biểu thức số
-------------[2] Biểu thức đại số
-------------[3] Giá trị của biểu thức số
-------------[4] Bài Toán Thực Tế
----------[2] Đa thức một biến
-------------[1] Đa thức một biến
-------------[2] Cách biểu diễn đa thức một biến
-------------[3] Giá trị của đa thức một biến
-------------[4] Nghiệm đa thức một biến
----------[3] Phép cộng và phép trừ đa thức một biến
-------------[1] Phép cộng hai đa thức một biến
-------------[2] Phép trừ hai đa thức một biến
-------------[3] Tính chất phép cộng đa thức một biến
-------------[4] Bài Toán Thực Tế
----------[4] Phép nhân và phép chia đa thức một biến
-------------[1] Phép nhân đa thức một biến
-------------[2] Phép chia đa thức một biến
-------------[3] Tính chất phép nhân đa thức một biến
-------------[4] Bài Toán Thực Tế
----[H] Hình học
-------[1] CÁC HÌNH KHỐI TRONG THỰC TIỄN
----------[1] Hình Hộp Chữ Nhật, Hình Lập Phương
-------------[1] Hình hộp chữ nhật
-------------[2] Hình lập phương
-------------[3] Bài Toán Thực Tế
----------[2] DTXQ VÀ THỂ TÍCH HHCN VÀ HÌNH LẬP PHƯƠNG
-------------[1] Nhắc lại công thức tính diện tích xung quanh và thể tích
-------------[2] Một số bài toán thực tế
----------[3] HÌNH LĂNG TRỤ ĐỨNG TAM GIÁC, HÌNH LĂNG TRỤ ĐỨNG TỨ GIÁC
-------------[1] Hình Lăng Trụ Đứng Tam giác, Hình Lăng Trụ Đứng Tứ Giác
-------------[2] Tạo lập Hình Lăng Trụ Đứng Tam giác, Hình Lăng Trụ Đứng Tứ Giác
-------------[3] Bài Toán Thực Tế
----------[4] DTQX VÀ THỂ TÍCH HÌNH LĂNG TRỤ ĐỨNG TAM GIÁC, TỨ GIÁC
-------------[1] Diện tích xung quanh hình lăng trụ đứng
-------------[2] Thể tích hình lăng trụ đứng
-------------[3] DTXQ và Thể Tích của một số hình khối trong thực tiễn
-------[2] HÌNH HỌC PHẲNG GÓC VÀ ĐƯỜNG THẰNG SONG SONG
----------[1] TÌM CÁC GÓC Ở VỊ TRÍ ĐẶC BIỆT
-------------[1] Hai góc kề bù
-------------[2] Hai góc đối đỉnh
-------------[3] Tính chất hai góc đối đỉnh
-------------[4] Bài Toán Thực Tế
----------[2] TIA PHÂN GIÁC
-------------[1] Tia Phân Giác Của Một Góc
-------------[2] Cách vẽ tia phân giác
-------------[3] Bài Toán Thực Tế
----------[3] Hai Đường Thẳng Song Song
-------------[1] Dấu hiệu nhận biết hai đường thẳng song song
-------------[2] Tiên đề Euclid về đường thẳng song song
-------------[3] Tính chất hai đường thẳng song song
-------------[4] Bài Toán Thực Tế
----------[4] Định Lí Và Chứng Minh Một Định Lí
-------------[1] Định Lí là gì?
-------------[2] Chứng Minh Định Lí
-------------[3] DTXQ và Thể Tích của một số hình khối trong thực tiễn
-------[3] HÌNH HỌC PHẲNG TAM GIÁC
----------[1] Góc Và Cạnh Của Một Tam Giác
-------------[1] Tổng số đo 3 góc của một tam giác
-------------[2] Quan hệ giữa 3 cạnh của một tam giác
-------------[3] Bài Toán Thực Tế
----------[2] Tam Giác Bằng Nhau
-------------[1] Hai tam giác bằng nhau
-------------[2] Trường hợp bằng nhau thứ nhất cạnh cạnh cạnh (c.c.c)
-------------[3] Trường hợp bằng nhau thứ hai cạnh góc cạnh (c.g.c)
-------------[4] Trường hợp bằng nhau thứ ba góc cạnh góc (g.c.g)
-------------[5] Trường hợp bằng nhau hai cạnh góc vuông (c.g.c)
-------------[6] Trường hợp bằng nhau cạnh góc vuông và một góc nhọn (g.c.g)
-------------[7] Trường hợp bằng nhau cạnh huyền và một góc nhọn (g.c.g)
-------------[8] Trường hợp bằng nhau cạnh huyền và một cạnh góc vuông (g.c.g)
-------------[9] Bài Toán Thực Tế
----------[3] Tam Giác Cân
-------------[1] Tam Giác Cân
-------------[2] Tính chất của tam giác cân
-------------[3] Bài Toán Thực Tế
----------[4] Đường Vuông Góc và Đường Xiên
-------------[1] Quan hệ giữa cạnh và góc trong một tam giác
-------------[2] Đường vuông góc và đường xiên
-------------[3] Mối quan hệ giữa đường vuông góc và đường xiên
-------------[4] Bài Toán Thực Tế
----------[5] Đường Trung Trực Của Một Đoạn Thẳng
-------------[1] Đường trung trực của một đoạn thẳng
-------------[2] Tính chất của đường trung trực
-------------[3] Bài Toán Thực Tế
----------[6] Tính Chất 3 Đường Trung Trực Của Tam Giác
-------------[1] Đường trung trực của tam giác
-------------[2] Tính chất 3 đường trung trực của tam giác
-------------[3] Bài Toán Thực Tế
----------[7] Tính Chất 3 Đường Trung Tuyến Của Tam Giác
-------------[1] Đường trung tuyến của tam giác
-------------[2] Tính chất 3 đường trung tuyến của tam giác
-------------[3] Bài Toán Thực Tế
----------[8] Tính Chất 3 Đường Cao Của Tam Giác
-------------[1] Đường cao của tam giác
-------------[2] Tính chất 3 đường cao của tam giác
-------------[3] Bài Toán Thực Tế
----------[9] Tính Chất 3 Đường Phân Giác Của Tam Giác
-------------[1] Đường phân giác của tam giác
-------------[2] Tính chất 3 đường phân giác của tam giác
-------------[3] Bài Toán Thực Tế
-[8] Lớp 8
----[D] Đại số
-------[1] Đa thức nhiều biến
----------[1] Đơn thức nhiều biến. Đa thức nhiều biến
-------------[1] Nhận dạng đơn thức
-------------[2] Nhận dạng đa thức
-------------[3] Thu gọn đơn thức
-------------[4] Tìm đơn thức đồng dạng
-------------[5] Cộng, trừ đơn thức đồng dạng
-------------[6] Thu gọn đa thức
-------------[7] Tính giá trị của đa thức
-------------[8] (Thêm) Bậc của đơn thức, đa thức
----------[2] Các phép tính với đa thức nhiều biến
-------------[1] Cộng hai đa thức nhiều biến
-------------[2] Trừ hai đa thức nhiều biến
-------------[3] Nhân hai đa thức nhiều biến
-------------[4] Chia hai đa thức nhiều biến
-------------[5] Rút gọn biểu thức
-------------[6] Chứng minh đẳng thức
-------------[7] Bài toán có lời văn
----------[3] Hằng đẳng thức đáng nhớ
-------------[1] Khai triển hằng đẳng thức
-------------[2] Viết biểu thức dưới dạng tích
-------------[3] Tính nhanh giá trị biểu thức
-------------[4] Chứng minh đẳng thức, biểu thức không phụ thuộc vào biến
----------[4] Vận dụng hằng đẳng thức vào phân tích đa thức thành nhân tử
-------------[1] Phân tích đa thức thành nhân tử bằng cách dùng hằng đẳng thức
-------------[2] Phân tích đa thức thành nhân tử bằng cách đặt nhân tử chung
-------------[3] Phân tích đa thức thành nhân tử bằng cách nhóm hạng tử
-------------[4] Phân tích đa thức thành nhân tử bằng cách kết hợp các phương pháp
-------------[5] Tính nhanh giá trị của biểu thức
-------------[6] Vận dụng phân tích đa thức thành nhân tử để chứng minh chia hết
-------------[7] Toán có lời văn
-------------[8] (Thêm) Phân tích đa thức thành nhân tử bằng cách tách hạng tử (đa thức bậc hai)
-------------[9] (Thêm) Phân tích đa thức thành nhân tử bằng cách thêm bớt hạng tử
-------------[0] Phân tích đa thức thành nhân tử bằng cách đặt ẩn phụ
-------[2] Phân thức đại số
----------[1] Phân thức đại số
-------------[1] Tìm điều kiện xác định của phân thức
-------------[2] Nhận dạng phân thức
-------------[3] Phân thức bằng nhau
-------------[4] Rút gọn phân thức
-------------[5] Quy đồng mẫu thức
-------------[6] Giá trị của một phân thức
-------------[7] Toán có lời văn
----------[2] Phép cộng, phép trừ phân thức đại số
-------------[1] Cộng hai hay nhiều phân thức
-------------[2] Tìm phân thức đối của một phân thức
-------------[3] Trừ hai hay nhiều phân thức
-------------[4] Rút gọn biểu thức có chứa phép cộng, trừ phân thức
-------------[5] Toán có lời văn
----------[3] Phép nhân, phép chia phân thức đại số
-------------[1] Nhân hai hay nhiều phân thức
-------------[2] Tìm phân thức nghịch đảo
-------------[3] Chia hai hay nhiều phân thức
-------------[4] Rút gọn biểu thức
-------------[5] Chứng minh đẳng thức, biểu thức không phụ thuộc vào biến
-------------[6] Toán có lời văn
-------[3] Hàm số và đồ thị
----------[1] Hàm số
-------------[1] Nhận dạng hàm số
-------------[2] Giá trị của hàm số
-------------[3] Toán có lời văn
----------[2] Mặt phẳng tọa độ. Đồ thị hàm số
-------------[1] Xác định tọa độ của điểm trên mặt phẳng tọa độ
-------------[2] Điểm thuộc (không thuộc) đồ thị hàm số
-------------[3] Toán có lời văn
----------[3] Hàm số bậc nhất y = ax + b (a khác 0)
-------------[1] Nhận dạng hàm số bậc nhất y = ax + b (a khác 0)
-------------[2] Xác định các yếu tố của hàm số bậc nhất y = ax + b (a khác 0)
-------------[3] Tính giá trị của hàm số
-------------[4] Toán có lời văn
----------[4] Đồ thị hàm số bậc nhất y = ax + b (a khác 0)
-------------[1] Xác định các yếu tố của hàm số bậc nhất
-------------[2] Vẽ đồ thị hàm số hàm số bậc nhất y = ax + b (a khác 0)
-------------[3] Tìm hệ số góc của đường thẳng
-------------[4] Xác định góc giữa đường thẳng và trục hoành
-------------[5] Vị trí tương đối của 2 đường thẳng
-------------[6] Xác định hàm số bậc nhất
-------------[7] Toán có lời văn
-------[4] Phương trình bậc nhất một ẩn
----------[1] Phương trình bậc nhất một ẩn
-------------[1] Nhận dạng phương trình
-------------[2] Nghiệm của phương trình
-------------[3] Nhận dạng phương trình bậc nhất một ẩn, xác định a, b
-------------[4] Giải phương trình bậc nhất một ẩn
-------------[5] Giải phương trình đưa được về dạng ax + b = 0
-------------[6] Toán có lời văn
----------[2] Ứng dụng của phương trình bậc nhất một ẩn
-------------[1] Bài toán về quan hệ giữa các số
-------------[2] Bài toán về chuyển động
-------------[3] Bài toán có nội dung hình học
-------------[4] Bài toán tỉ lệ, phần trăm
-------------[5] Bài toán về năng suất làm việc
-------[5] Một số yếu tố thống kê và xác suất
----------[1] Thu thập và phân loại dữ liệu
-------------[1] Thu thập dữ liệu
-------------[2] Phân loại và tổ chức dữ liệu
-------------[3] Tính hợp lý của dữ liệu
----------[2] Mô tả và biểu diễn dữ liệu trên các bảng, biểu đồ
-------------[1] Biểu diễn dữ liệu bằng bảng thống kê
-------------[2] Biểu diễn dữ liệu bằng biểu đồ cột
-------------[3] Biểu diễn dữ liệu bằng biểu đồ hình quạt
-------------[3] Biểu diễn dữ liệu bằng biểu đồ đoạn thẳng
-------------[5] Khai thác thông tin từ một tập dữ liệu
----------[3] Phân tích và xử lí dữ liệu thu được ở dạng bảng, biểu đồ
-------------[1] Phân tích và xử lí dữ liệu bằng bảng
-------------[2] Phân tích và xử lí dữ liệu bằng biểu đồ
-------------[3] Giải quyết những vấn đề đơn giản dựa trên phân tích và xử lí dữ liệu
----------[4] Xác suất của biến cố ngẫu nhiên trong một số trò chơi đơn giản
-------------[1] Xác suất của biến cố trong trò chơi tung đồng xu
-------------[2] Xác suất của biến cố trong trò chơi vòng quay số
-------------[3] Xác suất của biến cố trong trò chơi chọn ngẫu nhiên một đối tượng từ một nhóm đối tượng
----------[5] Xác suất thực nghiệm của một biến cố trong một số trò chơi đơn giản
-------------[1] Xác suất thực nghiệm của biến cố trong trò chơi tung đồng xu
-------------[2] Xác suất thực nghiệm của biến cố trong trò chơi gieo xúc xắc
-------------[3] Xác suất thực nghiệm của biến cố trong trò chơi chọn ngẫu nhiên một đối tượng từ một nhóm đối tượng
----[H] Hình học
-------[1] Hình học trực quan
----------[1] Hình chóp tam giác đều
-------------[1] Xác định đỉnh, cạnh, mặt hoặc nhận dạng hình chóp tam giác đều
-------------[2] Tính diện tích xung quanh của hình chóp tam giác đều
-------------[3] Xác định trung đoạn, chu vi đáy của hình chóp tam giác đều
-------------[4] Tính thể tích của hình chóp tam giác đều
-------------[5] Xác định đường cao, diện tích đáy của hình chóp tam giác đều
----------[2] Hình chóp tứ giác đều
-------------[1] Xác định đỉnh, cạnh, mặt hoặc nhận dạng hình chóp tứ giác đều
-------------[2] Tính diện tích xung quanh của hình chóp tứ giác đều
-------------[3] Xác định trung đoạn, chu vi đáy của hình chóp tứ giác đều
-------------[4] Tính thể tích của hình chóp tứ giác đều
-------------[5] Xác định đường cao, diện tích đáy của hình chóp tứ giác đều
-------[2] Tam giác. Tứ giác
----------[1] Định lý Pythagore
-------------[1] Tính độ dài cạnh của tam giác vuông
-------------[2] Chứng minh tam giác là tam giác vuông
-------------[3] Nhận dạng tam giác vuông dựa vào độ dài
-------------[4] Toán có lời văn
----------[2] Tứ giác
-------------[1] Nhận dạng tứ giác lồi
-------------[2] Tìm số đo góc của tứ giác
-------------[3] Toán có lời văn
----------[3] Hình thang cân
-------------[1] Nhận dạng hình thang cân
-------------[2] Chứng minh tứ giác là hình thang cân
-------------[3] Tính chu vi, diện tích, cạnh, đường cao của hình thang cân
-------------[4] Vận dụng tính chất của hình thang cân để chứng minh các vấn đề liên quan
-------------[5] Toán có nội dung thực tế
----------[4] Hình bình hành
-------------[1] Nhận dạng hình bình hành
-------------[2] Chứng minh tứ giác là hình bình hành
-------------[3] Tính chu vi, diện tích, cạnh, đường cao của hình bình hành
-------------[4] Vận dụng tính chất của hình bình hành để chứng minh các vấn đề liên quan
-------------[5] Toán có nội dung thực tế
----------[5] Hình chữ nhật
-------------[1] Nhận dạng hình chữ nhật
-------------[2] Chứng minh tứ giác là hình chữ nhật
-------------[3] Tính chu vi, diện tích, cạnh, đường cao của hình chữ nhật
-------------[4] Vận dụng tính chất của hình chữ nhật để chứng minh các vấn đề liên quan
-------------[5] Toán có nội dung thực tế
----------[6] Hình thoi
-------------[1] Nhận dạng hình thoi
-------------[2] Chứng minh tứ giác là hình thoi
-------------[3] Tính chu vi, diện tích, cạnh, đường cao của hình thoi
-------------[4] Vận dụng tính chất của hình thoi để chứng minh các vấn đề liên quan
-------------[5] Toán có nội dung thực tế
----------[7] Hình vuông
-------------[1] Nhận dạng hình vuông
-------------[2] Chứng minh tứ giác là hình vuông
-------------[3] Tính chu vi, diện tích, cạnh, đường cao của hình vuông
-------------[4] Vận dụng tính chất của hình vuông để chứng minh các vấn đề liên quan
-------------[5] Toán có nội dung thực tế
-------[3] Định lí Thales
----------[1] Định lí Thalès trong tam giác
-------------[1] Độ dài đoạn thẳng, đoạn thẳng tỉ lệ
-------------[2] Tính độ dài đoạn thẳng
-------------[3] Chứng minh hai đường thẳng song song
-------------[4] Chứng minh đẳng thức hình học
----------[2] Ứng dụng của định lí Thalès trong tam giác
-------------[1] Ước lượng khoảng cách
-------------[2] Ước lượng chiều cao
----------[3] Đường trung bình của tam giác
-------------[1] Tìm độ dài đường trung bình của tam giác
-------------[2] Tìm độ dài cạnh của tam giác
-------------[3] Vận dụng tính chất đường trung bình của tam giác để chứng minh các vấn đề liên quan
----------[4] Tính chất đường phân giác của tam giác
-------------[1] Tính độ dài đoạn thẳng
-------------[2] Chứng minh đẳng thức
-------------[3] Vận dụng tính chất đường phân giác của tam giác để chứng minh các vấn đề liên quan
-------[4] Tam giác đồng dạng. Hình đồng dạng
----------[5] Tam giác đồng dạng
-------------[1] Viết kí hiệu 2 tam giác đồng dạng
-------------[2] Tìm độ dài cạnh, góc khi biết 2 tam giác đồng dạng
-------------[3] Chỉ ra số tam giác đồng dạng
-------------[4] Chứng minh tam giác đồng dạng
----------[6] Trường hợp đồng dạng thứ nhất của tam giác
-------------[1] Chứng minh tam giác đồng dạng theo trường hợp c-c-c
-------------[2] Chứng minh tam giác vuông đồng dạng theo trường hợp ch-cgv
-------------[3] Vận dụng tam giác đồng dạng để chứng minh các vấn đề liên quan
-------------[4] Ứng dụng vào các bài toán thực tế
----------[7] Trường hợp đồng dạng thứ hai của tam giác
-------------[1] Chứng minh tam giác đồng dạng theo trường hợp c-g-c
-------------[2] Chứng minh tam giác vuông đồng dạng theo trường hợp 2 cạnh góc vuông
-------------[3] Vận dụng tam giác đồng dạng để chứng minh các vấn đề liên quan
-------------[4] Ứng dụng vào các bài toán thực tế
----------[8] Trường hợp đồng dạng thứ ba của tam giác
-------------[1] Chứng minh tam giác đồng dạng theo trường hợp g-g
-------------[2] Chứng minh tam giác vuông đồng dạng theo trường hợp góc nhọn
-------------[3] Vận dụng tam giác đồng dạng để chứng minh các vấn đề liên quan
-------------[4] Ứng dụng vào các bài toán thực tế
----------[9] Hình đồng dạng
-------------[1] Hình đồng dạng phối cảnh
-------------[2] Hình đồng dạng
-------------[3] Vận dụng hình đồng dạng để chứng minh hai hình bằng nhau
----------[0] Hình đồng dạng trong thực tiễn
-------------[1] Hình đồng dạng trong thế giới tự nhiên
-------------[2] Hình đồng dạng trong nghệ thuật, kiến trúc
-------------[3] Hình đồng dạng trong khoa học, công nghệ
----[G] Bồi dưỡng HSG
-------[1] Một số chủ đề Đại số, Số học bồi dưỡng HSG lớp 8
----------[1] Biểu thức đại số
----------[2] Phương trình, bất phương trình. Hệ phương trình, hệ bất phương trình
----------[3] Đa thức
----------[4] Dãy số nguyên
----------[5] Số học (tính chia hết, ƯCLN, BCNN, số nguyên tố, hợp số, số chính phương, phần nguyên,...)
----------[6] Phương trình nghiệm nguyên
----------[7] Đồng dư
----------[8] Bất đẳng thức, giá trị lớn nhất, giá trị nhỏ nhất
----------[9] Suy luận toán học; các bài toán về trò chơi; đại số tổ hơp
-------[2] Một số chủ đề Hình học bồi dưỡng HSG lớp 8
----------[1] Tứ giác và đường trung bình
----------[2] Diện tích đa giác
----------[3] Định lí Thales, tính chất phân giác
----------[4] Tam giác đồng dạng
----------[5] Các bài toán sử dụng các định lí hình học cổ điển (Ceva, Menelaus, ...)
----------[6] Chứng minh ba điểm thẳng hàng, ba đường thẳng đồng quy
----------[7] Yếu tố cố định
----------[8] Cực trị hình học
----------[9] Hình học tổ hợp
-[9] Lớp 9
----[D] Phần số
-------[1] PHƯƠNG TRÌNH VÀ HỆ PHƯƠNG TRÌNH
----------[1] PHƯƠNG TRÌNH QUY VỀ PHƯƠNG TRÌNH BẬC NHẤT
-------------[1] Phương trình tích
-------------[2] Phương trình chứa ẩn ở mẫu quy về phương trình bậc nhất
-------------[3] Giải bài toán bằng cách lập phương trình
----------[2] PHƯƠNG TRÌNH BẬC NHẤT 2 ẨN VÀ HỆ PHƯƠNG TRÌNH BẬC NHẤT 2 ẨN
-------------[1] Phương trình bậc nhất hai ẩn
-------------[2] Hệ phương trình bậc nhất hai ẩn
-------------[3] Bài Toán Thực Tế
----------[3] GIẢI HỆ PHƯƠNG TRÌNH BẬC NHẤT 2 ẨN
-------------[1] Giải hệ phương trình bằng phương pháp thế
-------------[2] Giải hệ phương trình bằng phương pháp cộng đại số
-------------[3] Tìm nghiệm của hệ bằng máy tính cầm tay
-------------[4] Giải bài toán bằng cách lập hệ phương trình
-------[2] BẤT ĐẲNG THỨC. BẤT PHƯƠNG TRÌNH BẬC NHẤT MỘT ẨN
----------[1] BẤT ĐẲNG THỨC
-------------[1] Khái niệm bất đẳng thức
-------------[2] Tính chất bất đẳng thức
-------------[3] Bài Toán Thực Tế
----------[2] BẤT PHƯƠNG TRÌNH BẬC NHẤT MỘT ẨN
-------------[1] Bất phương trình bậc nhất một ẩn và nghiệm của nó
-------------[2] Cách giải bất phương trình bậc nhất một ẩn
-------------[3] Bài Toán Thực Tế
-------[3] CĂN THỨC
----------[1] CĂN BẬC HAI
-------------[1] Căn bậc hai
-------------[2] Tính căn bậc hai bằng máy tính cầm tay
-------------[3] Căn thức bậc hai
-------------[4] Bài toán thực tế
----------[2] CĂN BẬC BA
-------------[1] Căn bậc ba của một số
-------------[2] Tính căn bậc ba bằng máy tính cầm tay
-------------[3] Căn thức bậc 3
-------------[4] Bài toán thực tế
----------[3] TÍNH CHẤT CỦA PHÉP KHAI PHƯƠNG
-------------[1] Căn thức bậc hai của một bình phương
-------------[2] Căn thức bậc hai của một tích
-------------[3] Căn thức bậc hai của một thương
-------------[4] Bài toán thực tế
----------[3] BIẾN ĐỔI ĐƠN GIẢN BIỂU THỨC CHỨA CĂN BẬC HAI
-------------[1] Trục căn thức ở mẫu
-------------[2] Rút gọn biểu thức chứa căn bậc hai
-------------[3] Bài toán thực tế
-------[4] HÀM SỐ y = ax^2 (a != 0) VÀ PHƯƠNG TRÌNH BẬC HAI MỘT ẨN
----------[1] HÀM SỐ VÀ ĐỒ THỊ HÀM SỐ y = ax^2 (a != 0)
-------------[1] Hàm số y = ax^2 (a != 0)
-------------[2] Bảng giá trị của hàm số y = ax^2 (a != 0)
-------------[3] Đồ thị của hàm số y = ax^2 (a != 0)
-------------[4] Bài toán thực tế
----------[2] PHƯƠNG TRÌNH BẬC 2 MỘT ẨN
-------------[1] Phương trình bậc hai một ẩn
-------------[2] Giải phương trình bậc hai bằng cách đưa về phương trình tích
-------------[3] Công thức nghiệm của phương trình bậc 2
-------------[4] Tìm nghiệm của phương trình bậc 2 bằng máy tính cầm tay
-------------[5] Giải bài toán bằng cách lập phương trình
----------[3] ĐỊNH LÝ VIÈTE
-------------[1] Định lí Viète
-------------[2] Tìm hai số khi biết tổng và tích của chúng
-------------[3] Bài toán thực tế
-------[5] MỘT SỐ YẾU THỐNG KÊ
----------[1] BẢNG TẦN SỐ VÀ BIỂU ĐỒ TẦN SỐ
-------------[1] Tần số và bảng tần số
-------------[2] Biểu đồ tần số
----------[2] BẢNG TẦN SỐ TƯƠNG ĐỐI VÀ BIỂU ĐỒ TẦN SỐ TƯƠNG ĐỐI
-------------[1] Bảng tần số tương đối
-------------[2] Biểu đồ tần số tương đối
----------[3] BIỂU DIỄN SỐ LIỆU GHÉP NHÓM
-------------[1] Bảng tần số ghép nhóm
-------------[2] Bảng tần số tương đối ghép nhóm
-------------[3] Biểu đồ tần số tương đối ghép nhóm
-------[6] MỘT SỐ YẾU TỐ XÁC SUẤT
----------[1] KHÔNG GIAN MẪU VÀ BIẾN CỐ
-------------[1] Không gian mẫu
-------------[2] Biến cố
----------[2] XÁC SUẤT CỦA BIẾN CỐ
-------------[1] Kết quả đồng khả năng
-------------[2] Xác suất của biến cố
----------[3] BIỂU DIỄN SỐ LIỆU GHÉP NHÓM
-------------[1] Bảng tần số ghép nhóm
-------------[2] Bảng tần số tương đối ghép nhóm
-------------[3] Biểu đồ tần số tương đối ghép nhóm
-------[7] MỘT SỐ BÀI TOÁN THỰC TẾ TRONG THI TUYỂN SINH
-------------[1] Bài Toán Can Chi
-------------[2] Bài Toán tính ngày này năm sau, năm trước là thứ mấy
-------------[3] Bài Toán tính múi giờ
-------------[4] Bài Toán tính chỉ số bằng phép cộng, trừ số nguyên
-------------[5] Bài Toán thiết lập đồ thị hàm số bậc nhất
-------------[6] Bài Toán thiết lập đồ thị hàm số bậc hai
-------------[7] Bài Toán ngân hàng, lãi suất
-------------[8] Bài Toán tối ưu (max min, ít nhất, nhiều nhất)
----[H] Hình học
-------[1] HỆ THỨC LƯỢNG TRONG TAM GIÁC VUÔNG
----------[1] TỈ SỐ LƯỢNG GIÁC CỦA GÓC NHỌN
-------------[1] Định nghĩa tỉ số lượng giác của góc nhọn
-------------[2] Tỉ số lượng giác của hai góc phụ nhau
-------------[3] Tính tỉ số lượng giác của góc nhọn bằng máy tính cầm tay
-------------[4] Bài toán thực tế
----------[2] HỆ THỨC GIỮA CẠNH VÀ GÓC CỦA TAM GIÁC VUÔNG
-------------[1] Hệ thức giữa cạnh và góc của tam giác vuông
-------------[2] Giải tam giác vuông
-------------[3] Bài toán thực tế
-------[2] ĐƯỜNG TRÒN
----------[1] ĐƯỜNG TRÒN
-------------[1] Khái niệm đường tròn
-------------[2] Tính chất đối xứng của đường tròn
-------------[3] Đường kính và dây cung
-------------[4] Vị trí tương đối của hai đường tròn
-------------[5] Bài toán thực tế
----------[2] TIẾP TUYẾN CỦA ĐƯỜNG TRÒN
-------------[1] Vị trí tương đối của đường thẳng và đường tròn
-------------[2] Dấu hiệu nhận biết tiếp tuyến của đường tròn
-------------[3] Tính chất hai tiếp tuyến cắt nhau
-------------[4] Bài toán thực tế
----------[3] GÓC Ở TÂM, GÓC NỘI TIẾP
-------------[1] Góc ở tâm
-------------[2] Cung, số đo cung
-------------[3] Góc nội tiếp
-------------[4] Bài toán thực tế
----------[4] HÌNH QUẠT TRÒN VÀ HÌNH VÀNH KHUYÊN
-------------[1] Độ dài cung tròn
-------------[2] Hình quạt tròn
-------------[3] Hình vành khuyên
-------------[4] Bài toán thực tế
-------[3] TỨ GIÁC NỘI TIẾP. ĐA GIÁC ĐỀU
----------[1] ĐƯỜNG TRÒN NGOẠI TIẾP TAM GIÁC. ĐƯỜNG TRÒN NỘI TIẾP TAM GIÁC
-------------[1] Đường tròn ngoại tiếp tam giác
-------------[2] Đường tròn nội tiếp tam giác
-------------[3] Bài toán thực tế
----------[2] TỨ GIÁC NỘI TIẾP
-------------[1] Định nghĩa tứ giác nội tiếp
-------------[2] Tính chất tứ giác nội tiếp
-------------[3] Đường tròn ngoại tiếp hình chữ nhật, hình vuông
-------------[4] Bài toán thực tế
----------[3] ĐA GIÁC ĐỀU VÀ PHÉP QUAY
-------------[1] Khái niệm đa giác đều
-------------[2] Phép quay
-------------[3] Hình phẳng đều trong thực tế
----------[4] BÀI HÌNH NHIỀU CÂU
-------------[1] Bài toán từ một điểm bên ngoài, kẻ hai tiếp tuyến đến (O)
-------------[2] Bài toán nửa đường tròn
-------------[3] Bài toán tam giác ABC có một cạnh là đường kính
-------------[4] Bài toán tam giác ABC nhọn nội tiếp đường tròn (O)
-------------[5] Bài toán dựng thêm đường tròn thứ hai khác (O)
-------------[6] Bài toán kẻ 2 dây cung vuông góc
-------------[7] Bài toán max-min
-------------[8] Bài toán khác
----------[5] BÀI HÌNH LIÊN QUAN CÁC HÌNH KHÁC
-------------[1] Bài toán kẻ 2 đường kính (hình chữ nhật)
-------------[2] Bài toán hình vuông
-------------[3] Bài toán hình tam giác đều
-------------[4] Bài toán hình tam giác cân
-------------[5] Bài toán hình tứ giác nội tiếp (cho 4 điểm thuộc đường tròn)
-------------[6] Bài toán hình thang
-------------[7] Bài toán hình bình hành
-------------[8] Bài toán max-min
----------[6] BÀI HÌNH LIÊN QUAN TIẾP TUYẾN CHUNG
-------------[1] Hai đường tròn tiếp xúc ngoài
-------------[2] Hai đường tròn tiếp xúc trong
-------------[3] Hai đường tròn ngoài nhau
-------------[4] Hai đường tròn đựng nhau
-------------[5] Hai đường tròn cắt nhau
-------[4] CÁC HÌNH KHỐI TRONG THỰC TIỄN
----------[1] HÌNH TRỤ
-------------[1] Hình trụ
-------------[2] Diện tích xung quanh hình trụ
-------------[3] Thể tích của hình trụ
-------------[4] Bài toán thực tế
----------[2] HÌNH NÓN
-------------[1] Hình nón
-------------[2] Diện tích xung quanh của hình nón
-------------[3] Thể tích của hình nón
-------------[4] Bài toán thực tế
----------[3] HÌNH CẦU
-------------[1] Hình cầu
-------------[2] Diện tích của mặt cầu
-------------[3] Thể tích của hình cầu
-------------[4] Bài toán thực tế
----------[4] CÁC HÌNH KHÔNG GIAN KẾT HỢP
-------------[1] Hình trụ và hình nón
-------------[2] Hình trụ và hình cầu
-------------[3] Hình cầu và hình nón
-------------[4] Các hình khác (lăng trụ, hình hộp)
-[A] Lớp 1
----[D] SỐ
-------[1] ALL
----------[1] ALL
-------------[1] ALL
----[H] HÌNH
-------[1] ALL
----------[1] ALL
-------------[1] ALL
-[B] Lớp 2
----[D] SỐ
-------[1] ALL
----------[1] ALL
-------------[1] ALL
----[H] HÌNH
-------[1] ALL
----------[1] ALL
-------------[1] ALL
-[C] Lớp 3
----[D] SỐ
-------[1] ALL
----------[1] ALL
-------------[1] ALL
----[H] HÌNH
-------[1] ALL
----------[1] ALL
-------------[1] ALL
-[D] Lớp 4
----[D] SỐ
-------[1] ALL
----------[1] ALL
-------------[1] ALL
----[H] HÌNH
-------[1] ALL
----------[1] ALL
-------------[1] ALL
-[E] Lớp 5
----[D] SỐ
-------[1] ALL
----------[1] ALL
-------------[1] ALL
----[H] HÌNH
-------[1] ALL
----------[1] ALL
-------------[1] ALL
%Kết thúc nội dung ID